% ---
% title: 用 LaTeX 写博客的第一篇文章
% date: 2026-07-01
% author: Crystal-Sky
% tags: LaTeX, Blog, MathJax
% summary: 这是一篇示例文章，展示如何在新的个人博客里使用 LaTeX 结构、数学公式、定理和证明。
% course: 学习与写作工具
% course_slug: learning-toolkit
% lesson_order: 1
% lesson_type: 入门指南
% ---

\section{为什么要独立建站}

博客园适合快速发布 Markdown 笔记，但当文章里有大量公式、定理环境、推导过程和需要长期维护的知识结构时，使用 \LaTeX 写作会更自然。这个站点把文章源文件放在 `posts/*.tex`，构建时生成静态 HTML，再交给 GitHub Pages 或其他静态服务部署。

\begin{abstract}
这个示例展示了文章元信息、章节、行内公式、展示公式、定理环境、列表和代码块。
\end{abstract}

\section{数学公式}

行内公式可以写成 $e^{i\pi}+1=0$，展示公式可以写成：

\[
  \sum_{k=1}^{n} k = \frac{n(n+1)}{2}.
\]

也可以使用 `align` 环境：

\begin{align}
  (a+b)^2 &= a^2 + 2ab + b^2,\\
  (a-b)^2 &= a^2 - 2ab + b^2.
\end{align}

\section{定理与证明}

\begin{theorem}[等式的平移]
若 $a=b$，则对任意 $c$，都有 $a+c=b+c$。
\end{theorem}

\begin{proof}
根据等式的基本性质，在等式两边同时加上同一个量 $c$，等式仍然成立。
\end{proof}

\section{列表与代码}

\begin{itemize}
\item 在 `posts` 目录中写 `.tex` 文章。
\item 运行 `npm run build` 生成静态网站。
\item 推送到 GitHub 后由 Actions 自动部署。
\end{itemize}

\begin{verbatim}
npm run new "新的文章标题"
npm run dev
npm run build
\end{verbatim}
